\section{实验目的}

本实验使用 OpenCV 官方样例数据中的 \texttt{aloeL.jpg} 和 \texttt{aloeR.jpg}，完成从左右图像到
视差图和相对深度图的完整处理流程。实验重点不是重新实现立体匹配算法，而是把
输入校验、预处理、匹配、数值筛选、结果保存和证据核对连接成一个可复现的实验闭环。

\subsection{课程意义}

双目深度估计把计算机视觉中的几何模型和实际图像处理接口联系起来。视差是两幅图像
之间可直接观测的像素位移，而深度还受到焦距、基线和标定质量的影响。通过本实验，
可以区分“得到一张视觉上有层次的图”与“得到具有物理单位、可用于测量的深度”这两个
不同目标，并理解无效匹配、遮挡和纹理条件对结果的影响。

\subsection{实验目标}

\begin{enumerate}
  \item 理解针孔双目模型、对极约束、视差与深度之间的关系；
  \item 掌握 OpenCV 图像读取、灰度预处理和 \texttt{StereoSGBM} 立体匹配接口；
  \item 将匹配器输出的定点视差转换为像素单位的浮点视差，并进行有效区域筛选与可视化；
  \item 解释 \texttt{minDisparity}、\texttt{numDisparities}、平滑项和 speckle 过滤等参数对匹配结果的作用；
  \item 在缺少与图像对配套的相机标定参数时，计算并分析无物理单位的相对深度代理；
  \item 保存真实运行产生的左右图、视差图、有效掩码、数组和参数摘要，并据此完成结果核验。
\end{enumerate}

\subsection{可复现性与输入输出边界}

本报告以 OpenCV 官方 \texttt{4.x/samples/data} 中的 aloe 图像对为唯一主实验输入。程序命令、
软件版本、匹配参数、输出数组尺寸和统计值均以本次运行生成的
\codepath{task2/results/params.json} 与 \codepath{run_evidence.md} 为准。实验输出包括
像素视差和任意尺度的
\(\mathrm{depth\_proxy}=1/d\)；由于当前图像对没有经过核验的焦距、基线和 $\mathbf{Q}$，
不把该代理解释为米、厘米等物理单位的绝对深度。
